\section*{AI使用报告}

\subsection*{1. AI 工具名称、版本或型号}

本参赛队在竞赛过程中使用的生成式人工智能工具如下。

\begin{center}
\begin{tabularx}{0.92\textwidth}{>{\centering\arraybackslash}p{2.6cm}>{\centering\arraybackslash}p{2.6cm}>{\centering\arraybackslash}p{3.0cm}X}
\toprule
\textbf{AI 工具} & \textbf{开发机构} & \textbf{使用模型} & \textbf{主要用途}\\
\midrule
ChatGPT & OpenAI & GPT-5.6 Sol & 文献阅读辅助、建模思路讨论、代码实现指导\\
\bottomrule
\end{tabularx}
\end{center}

AI 工具主要作为资料检索与理解、候选数学表达讨论和程序实现辅助工具。题目分析、模型取舍、模型修正、参数解释、计算结果判断及论文最终内容均由参赛队成员完成。

\subsection*{2. 具体使用目的和环节}

\noindent\textbf{1. 文献阅读辅助}\par
在确定防护服多层瞬态传热建模、DSC相变数据处理及拓展分离变量法求解方案时，使用 AI 辅助阅读相关论文。主要用途包括：
\begin{itemize}
  \item 提取相变材料DSC曲线处理、操作基线选择和潜热反演方法；
  \item 比较不同层合材料瞬态传热求解方法（拓展分离变量法、有限体积法等）的适用条件；
  \item 根据文献中的数学模型、边界条件设置和数值求解策略进行人工检索和复现。
\end{itemize}

\noindent\textbf{2. 建模思路讨论}\par
在四个问题的模型形成过程中，使用 AI 对候选数学问题和建模方法进行讨论，主要涉及：
\begin{itemize}
  \item 人体恒温边界下三层防护服耦合瞬态传热的数学模型建立；
  \item 基于DSC数据的有效比热构造与不可逆固化进度处理方式；
  \item 外对流换热系数在自然对流与强制对流条件下的确定方法；
  \item 离散厚度优化中热安全约束与负重约束的耦合建模；
  \item 相变材料放热能力同比例增强的反问题建模与二分搜索求解。
\end{itemize}

\noindent\textbf{3. 代码实现指导}\par
在模型结构由参赛队确定后，使用 AI 辅助讨论部分数值算法的实现，包括：
\begin{itemize}
  \item 微小时间段拓展分离变量法的数值流程与界面斜率闭合法则；
  \item 基于DSC校正峰的有效比热插值与不可逆相变历史约束更新；
  \item 二分搜索框架下功能层放热能力倍率与有效持续时间的联动计算；
  \item 离散枚举方案中热安全时间与负重时间的比较与最优筛选。
\end{itemize}

\subsection*{3. 主要提示方式与使用过程}

\noindent\textbf{Example 1. OpenAI ChatGPT（GPT-5.6 Sol）— 背景与建模讨论}\par

\textbf{Query:} 〈先调研人体在低温环境下的防护服瞬态传热建模方法。为什么相变材料在防护服中能够延缓人体表面温度下降？DSC数据中负峰如何解释为放热过程？操作基线选择对反演的绝对潜热值有多大影响？先解释这些物理和数据处理问题，再判断它们在题目中应如何纳入数学模型。〉

\textbf{Output:} 〈相变材料在温度降至凝固点附近时释放固化潜热，该放热过程在DSC曲线上表现为负峰。对DSC数据取反并构造基线后，可由峰面积与扫描速率的比值计算单位质量潜热。操作基线选择（如端点连线、切线基线等）会直接影响积分面积，进而改变反演的潜热值。

在本题背景下，可将DSC校正曲线转换为固化进度随温度的分布，再通过有效比热形式将相变效应写入第二层的瞬态导热方程中，并采用历史最大值约束保证相变不可逆。人体温度按附件测试条件保持37°C恒定，作为恒温边界处理。〉

\noindent\textbf{Example 2. OpenAI ChatGPT（GPT-5.6 Sol）— 代码实现指导}\par

\textbf{Query:} 〈我们准备按照拓展分离变量法求解三层防护服耦合瞬态传热模型。先对DSC数据做基线校正并构造固化进度，再在每个微小时间段内将两处界面温度线性外推，通过本征函数展开推进温度场。如何设计这个程序？〉

\textbf{Output:} 〈程序可分为五个核心模块：第一，读取DSC数据并完成基线校正，由校正峰面积和扫描速率假设计算潜热，构造固化进度与有效比热插值函数；第二，预计算三层各自的本征值与本征函数，第一层采用Robin-Dirichlet、第二层采用Dirichlet-Dirichlet、第三层采用Dirichlet-Robin边界；第三，初始化温度场和不可逆固化进度，在每个时间步内由当前功能层平均温度更新有效比热；第四，利用两处界面温度线性外推构造非齐次项的模态投影，由段末热流连续条件联立求解界面斜率；第五，推进模态系数后更新温度场，并在阈值穿越时刻插值输出。需要额外检查功能层平均温度是否触发新的相变进度、固化完成后有效比热退回显热比热以及模态截断数增加后的收敛性。〉

\subsection*{4. AI 输出的采纳、人工修改和核验情况}

\noindent\textbf{1. DSC基线选择的人工判断}\par
AI 曾建议采用切线基线或多项式基线处理DSC放热峰。参赛队指出，题目未给出固化的明确起止边界和基线形态信息，附件数据仅包含温度-热流曲线；复杂的基线选择会引入额外假设且难以验证。因此采用温区两端点连线作为操作基线，该规则直接利用数据自身端点，不引入可调参数，且可复现。

\noindent\textbf{2. 外对流相关式选择经文献核实修正}\par
AI 根据热舒适与冷应激文献建议了ISO 11079边界空气层热阻关系和PMV对照关系。参赛队核对文献后确认，ISO 11079给出的热阻关系适用于冷环境风速条件，但原文按有效面积和辐射项给出完整形式，不能直接截取单一系数；直接使用该关系时须做等有效面积简化且忽略辐射项。因此正式模型仅采用其边界空气层热阻核心关系，并明确记录简化边界；PMV和Kuwabara关系仅用于敏感性对照，不作为主结果。

\noindent\textbf{3. 数值方法与参数假设均经人工判断}\par
AI 给出的DSC扫描速率假设（基准取5 K/min）及敏感性范围（2和10 K/min）均不作为最终结论本身。参赛队通过正式程序重新计算，并通过以下方式进行核验：
\begin{itemize}
  \item 检查DSC校正峰面积与潜热的转换是否满足量纲一致性；
  \item 比较基准工况下模态法结果与独立隐式有限体积法的相对偏差；
  \item 对时间步减半和模态数加密检查阈值时间是否达到数值收敛；
  \item 对DSC扫描速率和衣料初温分别做敏感性分析，评估建模假设对结果的影响幅度；
  \item 对问题三的低扫描速率诊断结果和问题四的二分搜索上下界收敛性做独立复核。
\end{itemize}

因此 AI 输出主要用于形成候选方案和减少程序实现中的重复工作。正式论文中的模型结构、适用条件、数值结果和结论均经过重新判断和验证。
